\section{Worked Examples}
\label{sec:appendix}

\lstset{basicstyle=\scriptsize\ttfamily, frame=single, captionpos=b,
        columns=fullflexible, keepspaces=true, breaklines=true,
        showstringspaces=false}

Three examples, each shown as its input, its state in between and its output,
cut down to what the example is meant to show. The complete files, and every
other example, are in the hand-in archive (Appendix~\ref{sec:handin}).

\subsection{The Minimal Model, End to End}
\label{sec:appendix-minimal}

The input is one project file in the concrete syntax of
Section~\ref{sec:constraints}. Comments are removed here; the committed file carries
one per decision.

\begin{lstlisting}[caption={Input: \texttt{minimal/notes.project.yml}}, label={lst:minimal-input}]
apiVersion: intent.jorisjonkers.dev/v1
kind: Project
schemaVersion: 1.0.0
project: notes
owner: joris
applications:
  - id: notes
    observability:
      alertClass: business-hours
      scrape: { process: notes-api, surface: http, path: /metrics }
    exposure:
      - name: public
        host: notes.jorisjonkers.dev
        audience: anonymous
        contentPolicy: strict
        routes:
          - { path: /, match: prefix, process: notes-api, surface: http }
    processes:
      - name: notes-api
        lifecycle: application
        image: notes-api
        runtime: node
        provides: { http: 8080 }
        placement: { memory: 256Mi, cpu: 50m }
        probes:
          readiness: { path: /healthz/ready, port: 8080 }
          liveness: { path: /healthz/live, port: 8080 }
        startupBudget: 20s
        cutover: continuous
\end{lstlisting}

The parser turns it into an instance of the source metamodel. The same instance,
serialised as XMI, is what Eclipse opens without the grammar; attributes that
hold their default value, such as \texttt{lifecycle: application}, are not
written.

\begin{lstlisting}[caption={Input as a model: \texttt{minimal/intent.xmi}}, label={lst:minimal-xmi}]
<projectintent:Project apiVersion="intent.jorisjonkers.dev/v1" kind="Project"
    schemaVersion="1.0.0" project="notes" owner="joris">
  <applications id="notes">
    <observability>
      <scrape process="notes-api" surface="http" path="/metrics"/>
    </observability>
    <exposure name="public" host="notes.jorisjonkers.dev" audience="anonymous" contentPolicy="strict">
      <routes path="/" process="notes-api" surface="http"/>
    </exposure>
    <processes startupBudget="20s" cutover="continuous" name="notes-api" image="notes-api" runtime="node">
      <placement memory="256Mi" cpu="50m"/>
      <provides key="http" value="8080"/>
      <probes xsi:type="projectintent:Probes">
        <readiness xsi:type="projectintent:HttpProbe" path="/healthz/ready" port="8080"/>
        <liveness xsi:type="projectintent:HttpProbe" path="/healthz/live" port="8080"/>
      </probes>
    </processes>
  </applications>
</projectintent:Project>
\end{lstlisting}

In between is the resolved model: an instance of the target metamodel, written
by hand as the result the Task 2 transformation must produce. Every value in it
is either carried from the input or decided by the platform, and nothing in it
names a Kubernetes field.

\begin{lstlisting}[caption={In between: \texttt{minimal.resolveddeployment}, the application}, label={lst:minimal-resolved}]
<applications id="notes" revision="sha256:..." project="notes"
    namespace="notes-system" alertClass="business-hours" reconcileUnit="//@reconcileUnits.0">
  <releaseGate deadline="60s">
    <analysis interval="30s" iterations="4" threshold="3"/>
    <members process="//@applications.0/@processes.0">
      <readiness xsi:type="resolveddeployment:HttpTarget" path="/healthz/ready" port="8080"/>
      <checks>error-rate</checks>
      <checks>latency</checks>
    </members>
  </releaseGate>
  <exposure name="public" host="notes.jorisjonkers.dev" tier="public-frankfurt">
    <routes path="/" match="prefix" audience="anonymous" precedence="1"
        process="//@applications.0/@processes.0" surface="http">
      <middleware kind="security-headers" contentPolicy="strict"/>
    </routes>
  </exposure>
  <processes name="notes-api" identity="notes-api" image="ghcr.io/.../notes-api@sha256:9c1e..."
      uid="1000" gid="1000" cutover="continuous" switchover="blue-green" deadline="60s"
      replicas="1" memory="256Mi" cpu="50m" hardening="restricted" identityToken="false">
    <startup period="5s" failures="4">
      <target xsi:type="resolveddeployment:HttpTarget" path="/healthz/live" port="8080"/>
    </startup>
    <placement eligibleNodes="enschede-t1000-1 ... frankfurt-contabo-1"/>
  </processes>
</applications>
\end{lstlisting}

Table~\ref{tab:minimal-trace} reads the two side by side: what the author wrote,
and what the platform decided from it.

\begin{table}[!htbp]
\centering
\footnotesize
\begin{tabular}{>{\raggedright\arraybackslash}p{0.30\linewidth}>{\raggedright\arraybackslash}p{0.62\linewidth}}
\hline
\textbf{Authored} & \textbf{Decided in the resolved model} \\
\hline
\texttt{project: notes} & \texttt{namespace="notes-system"}, and the reconcile unit \\
\texttt{cutover: continuous} & \texttt{switchover="blue-green"}, and a \texttt{releaseGate} with this process as its one member \\
\texttt{startupBudget: 20s} & a startup probe of 4 checks 5 seconds apart, and a deadline of 60 seconds \\
\texttt{runtime: node} & the checks \texttt{error-rate} and \texttt{latency}, because the profile exposes HTTP server metrics \\
\texttt{image: notes-api} & the image digest, and the numeric user and group it runs as \\
\texttt{audience: anonymous} & the edge tier \texttt{public-frankfurt}, and a middleware chain without forward authentication \\
\texttt{placement} & the machines the process is eligible for: all seven, because nothing narrows the set \\
\hline
\end{tabular}
\caption{The Minimal Model: What Was Authored, and What Was Decided}
\label{tab:minimal-trace}
\end{table}

The output is the generated file tree: nine files for this one process. The
Canary is the object through which the process is switched; every decision in it
was already in the resolved model, and only the generator spells it.

\begin{lstlisting}[caption={Output: \texttt{minimal/rendered/apps/notes/canary.yaml}, one webhook of three}, label={lst:minimal-output}]
apiVersion: flagger.app/v1beta1
kind: Canary
metadata:
  name: notes-api
  namespace: notes-system
spec:
  provider: kubernetes
  targetRef: { apiVersion: apps/v1, kind: Deployment, name: notes-api }
  progressDeadlineSeconds: 60
  service: { port: 8080, targetPort: http, portName: http }
  analysis:
    interval: 30s
    iterations: 4
    threshold: 3
    webhooks:
    - name: may-promote
      type: confirm-promotion
      url: http://release-gate.delivery-system.svc.cluster.local:8080/may-promote
      metadata: { application: notes, process: notes-api, revision: "sha256:..." }
\end{lstlisting}

\subsection{Two Cutovers in One Project}
\label{sec:appendix-knowledge}

The \texttt{knowledge} example shows how one authored question changes the
shape of the output. Its two applications answer the cutover question
differently, so they are two applications, and only one of them moves the
project's schema.

\begin{lstlisting}[caption={Input: \texttt{knowledge/knowledge.project.yml}, abridged}, label={lst:knowledge-input}]
applications:
  - id: knowledge
    migration:
      changelog: knowledge-api/src/main/resources/db/changelog.yml
    processes:
      - name: knowledge-api
        runtime: jvm
        startupBudget: 600s
        cutover: continuous
  - id: knowledge-ingest
    migration: none
    processes:
      - name: knowledge-ingest-worker
        runtime: python
        probes: none
        cutover: interrupted
        volumes:
          - { claim: knowledge-vault-clone, size: 20Gi, durability: irreplaceable }
\end{lstlisting}

In between, the continuous application carries a release gate and a migration;
the interrupted one carries neither, because it stops before it starts and
nothing waits on a gate.

\begin{lstlisting}[caption={In between: the two applications' projections, abridged}, label={lst:knowledge-resolved}]
knowledge:
  "migration":   { "runner": ".../knowledge-migration@sha256:c0a4...",
                   "testedAgainst": "sha256:5a1f...", "nonTransactional": false },
  "releaseGate": { "deadline": "1800s",
                   "members": [ { "process": "knowledge-api",
                                  "checks": ["error-rate", "latency"] } ] },
  "processes":   [ { "name": "knowledge-api", "switchover": "blue-green" } ]

knowledge-ingest:
  "processes":   [ { "name": "knowledge-ingest-worker", "cutover": "interrupted",
                     "switchover": "stop-start" } ]
\end{lstlisting}

In the output, the migration becomes two Jobs named for the revision: the one
that moves the schema, and the one that would undo it, rendered suspended so
that only the release gate can start it.

\begin{lstlisting}[caption={Output: \texttt{knowledge/rendered/apps/knowledge/migration.yaml}, the down, abridged}, label={lst:knowledge-output}]
apiVersion: batch/v1
kind: Job
metadata:
  name: knowledge-migration-down-6325500a2c54
  annotations:
    kustomize.toolkit.fluxcd.io/ssa: IfNotPresent
spec:
  suspend: true
  backoffLimit: 0
  template:
    spec:
      serviceAccountName: knowledge-migration
      containers:
      - name: migration
        image: ghcr.io/jorisjonkers-dev/knowledge/knowledge-migration@sha256:c0a4...
        args: [down, --to=sha256:5a1f...]
\end{lstlisting}

\subsection{A Refused Model}
\label{sec:appendix-refused}

A prepare process runs once before the new version starts, so it serves
nothing. Declaring what only a serving process has is refused, and the
diagnostic names the object at fault by its JSON Pointer.

\begin{lstlisting}[caption={Input: \texttt{refusals/prepare-process-serves.project.yml}, the refused process}, label={lst:refused-input}]
      - name: seed
        lifecycle: prepare
        image: seed
        runtime: none
        placement: { memory: 64Mi, cpu: 10m }
        startupBudget: 60s
        provides: { http: 8080 }
        probes:
          readiness: { path: /ready, port: 8080 }
        replicas: { count: 2, reason: A step that runs once has no second copy to spread. }
        cutover: continuous
\end{lstlisting}

\begin{lstlisting}[caption={Output: \texttt{refusals/prepare-process-serves.diagnostics.json}}, label={lst:refused-output}]
[{"code":"E_PREPARE_PROCESS_SERVES","document":"prepare-process-serves.project.yml",
  "path":"/applications/0/processes/1"}]
\end{lstlisting}

\subsection{The Hand-in Archive}
\label{sec:handin}

The archive handed in with this report holds the metamodels and every model
referred to here. Table~\ref{tab:handin} lists its layout. Each example has
the same three stages as the examples above. The resolved model exists as XMI
for the minimal example and as the production implementation's JSON projection
for \texttt{minimal}, \texttt{knowledge} and \texttt{delivery}; the others have
none yet, because producing it is the Task 2 transformation.

\begin{table}[!htbp]
\centering
\footnotesize
\begin{tabular}{>{\raggedright\arraybackslash}p{0.36\linewidth}>{\raggedright\arraybackslash}p{0.56\linewidth}}
\hline
\textbf{Path} & \textbf{What it holds} \\
\hline
\texttt{metamodels/} & \texttt{project-intent.ecore}, \texttt{resolved-deployment.ecore}, and \texttt{project-intent.ocl}, the constraints of Section~\ref{sec:constraints} \\
\texttt{examples/<name>/1-input/} & the authored files in concrete syntax, and \texttt{intent.xmi}, the same model as an instance of the source metamodel \\
\texttt{examples/<name>/2-resolved/} & the resolved model: \texttt{minimal.resolveddeployment} as XMI, and \texttt{resolved*.json} projections \\
\texttt{examples/<name>/3-output/} & the generated file tree \\
\texttt{examples/platform/} & the Platform document, and its model as XMI \\
\texttt{refusals/} & every refused model with the diagnostics it must produce \\
\hline
\end{tabular}
\caption{The Layout of the Hand-in Archive}
\label{tab:handin}
\end{table}
